\section{Introduction}
\label{sec:intro}

Automated software testing plays a pivotal role in modern software quality assurance, enabling continuous integration pipelines to validate system behavior across complex codebases~\cite{myers2011art, cse}. Comprehensive unit testing is essential to prevent subtle code regressions, enforce architectural invariants, and validate intended system semantics across interdependent modules~\cite{docprompting, repobench}. Software testing is a key component of the engineering workflows at major technology enterprises such as Google, Microsoft, and Meta~\cite{swebench}. However, writing and maintaining test suites manually remains a major cost driver in software development: in industrial settings, testing activities account for up to 50\% of a project's total development budget~\cite{myers2011art, bareiss2022code}.

The importance of software testing has driven the creation of numerous techniques and tools for the automated generation of test cases~\cite{evosuite, randoop, schafer2023empirical}. Recent advancements in Large Language Models (LLMs) have enabled generative tools to synthesize readable unit tests and achieve high code coverage across diverse programming languages~\cite{bareiss2022code, lemieux2023codamosa, xie2023chatunitest}. Most existing LLM-based test generators adopt a static, single-pass completion approach, where test assertions are derived directly from code signatures and surrounding context in a "Plain Context" paradigm~\cite{testexplora}. However, these static tools are fundamentally limited in their ability to detect logic flaws because they suffer from \textit{compliance bias}: when presented with a defective repository implementation, static models passively generate test assertions that conform to the existing (buggy) implementation rather than validating intended behavior~\cite{testexplora, sweagent}. For instance, when generating a test for a defective calculation function that incorrectly returns zero on negative inputs, a static LLM will generate \texttt{assert calc(-5) == 0}, causing the test suite to pass on the buggy codebase while failing to expose the underlying defect. None of these logic errors would be detected by existing static completion tools since all generated assertions syntactically conform to the defective code. Recent benchmarks~\cite{testexplora} have identified breaking compliance bias in repository-level test generation as a major open challenge in automated software engineering. This is the problem that motivates our work.

The automated generation of test suites is an active research topic. Existing approaches mostly differ on the inputs and signals used for test generation, including source code signatures~\cite{bareiss2022code}, search-based genetic mutation~\cite{evosuite, randoop}, and documentation-derived intent~\cite{docprompting, testexplora}. One promising direction is the integration of interactive execution feedback, allowing testing agents to observe runtime tracebacks and iteratively refine test assertions. This is the strategy we use.

This article presents an open-source, dual-agent framework for proactive repository-level test suite generation through interactive terminal execution feedback. Our approach pairs an Exploration Agent (DeepSeek-V3) operating through an interactive bash terminal via the SWE-agent CLI with a Code Action Fixer Agent (Llama-3.3-70B). By executing candidate test suites in real-time, analyzing \texttt{pytest} tracebacks, and iteratively refining adversarial assertions across up to 80 interaction turns, our system dynamically rejects compliance assumptions and isolates repository-level logic defects without relying on human-written issue reports.

Evaluation results using $N=1,552$ tasks across 482 public Python repositories from the TestExplora benchmark showed that interactive terminal execution feedback significantly outperforms static baseline generators. At a single-sample budget ($k=1$), our approach achieved an average Fail-to-Pass (F2P) rate of 35.05\%, more than double the static baseline's 16.60\%. When expanding the sample budget to $K=3$ independent runs (Pass@3), the cumulative bug discovery rate reached 71.59\% ($p = 0.000000$, Cliff's $\delta = 0.3018$, medium effect size). During our evaluation, interactive execution tracebacks allowed the agent to expose latent logic bugs across 10 distinct software domain categories---including database query engines, infrastructure monitoring tools, and web frameworks---that passed completely unnoticed by static completion tools.

This article first introduces background and related work on automated test generation, LLM compliance bias, and execution-driven agentic interfaces (Section~\ref{sec:related}). Then, the article presents the following original research and engineering contributions:
\begin{itemize}
    \item \textbf{Dual-Agent Terminal Framework:} An open-source multi-agent system combining DeepSeek-V3 via SWE-agent CLI bash terminal with Llama-3.3-70B for proactive repository-level test generation (Section~\ref{sec:method}).
    \item \textbf{Fault-Tolerant Parallel Infrastructure:} A 16-shard round-robin distribution pipeline equipped with exponential backoff retries and execution timeouts to process large-scale agentic runs reliably (Section~\ref{subsec:setup}).
    \item \textbf{Comprehensive Empirical Evaluation:} Large-scale experimental results on $N=1,552$ tasks ($4,656$ total agent runs) demonstrating a 35.05\% Pass@1 and 71.59\% Pass@3 F2P rate, supported by non-parametric statistical hypothesis testing ($p = 0.000000$, Cliff's $\delta = 0.3018$) (Section~\ref{sec:results}).
    \item \textbf{Domain \& Failure Mode Analysis:} Qualitative investigation of 10 software domain categories and empirical breakdown of unresolved edge cases (Section~\ref{sec:discussion}).
    \item \textbf{Open Replication Package:} A publicly available dataset, execution logs, analysis scripts, and figure generation pipeline hosted on GitHub at \url{https://github.com/lctx9/group4} (Section~\ref{subsec:setup}).
\end{itemize}

Section~\ref{sec:threats} addresses threats to internal, external, construct, and conclusion validity, and Section~\ref{sec:conclusion} concludes the paper and outlines future work.
